\section{Proofs for Section~\ref{sec:theory}}\label{app:theory}

\paragraph{Proposition~\ref{thm:lineage}.}
\emph{(a)} Let $N_k$ be the event that exam $k$ does not admit. Given $\mathcal H_{K-1}$ on the event $N_1\cap\dots\cap N_{K-1}$, exam $K$ is valid at level $a_K$ by Definition~\ref{def:condnull}. Hence $\Pr(N_1\cap\dots\cap N_K)=\mathbb E[\mathbf 1\{N_1\cap\dots\cap N_{K-1}\}\Pr(N_K\mid\mathcal H_{K-1})]\ge(1-a_K)\Pr(N_1\cap\dots\cap N_{K-1})$. Induction gives $\Pr(\bigcap_{k\le K}N_k)\ge\prod_{k\le K}(1-a_k)$; let $K\to\infty$, and use $1-\prod_k(1-a_k)\le\sum_ka_k$.
\emph{(b)} Each exam's rule is fixed in advance, its variants are drawn independently of earlier exams, and the version's score law $Q$ does not change, so the outcome of exam $k$ is a fixed function of its own variants and the events $N_k$ are independent with $\Pr(N_k)=1-\sigma_k(Q)$, and $\Pr(\text{ever admitted})=1-\prod_k(1-\sigma_k(Q))$. Control at $\alpha$ requires $\prod_k(1-\sigma_k(Q))\ge1-\alpha$, that is, $\sum_k-\log(1-\sigma_k(Q))\le\bar\alpha$, and $-\log(1-s)\ge s$. With $\sigma_k(Q)\ge\sigma$ for all $k$ the error after $K$ attempts is at least $1-(1-\sigma)^K$.

\paragraph{Proposition~\ref{thm:onexam}.}
\emph{(a)} If $\mathbb E_P[\tau]=\infty$ there is nothing to prove. Take any null $Q$ with $\mathrm{KL}(P\|Q)<\infty$ and let $A$ be the admission event. Under $Q$ the exam may never stop, so we work with the bounded stopping time $\tau_n=\tau\wedge n$ and the event $A_n=A\cap\{\tau\le n\}$, which is $\mathcal F_{\tau_n}$-measurable. The change-of-measure inequality \citep[Lemma~1]{kaufmann2016complexity} gives $\mathbb E_P[\tau_n]\,\mathrm{KL}(P\|Q)\ge\mathrm{kl}(P(A_n),Q(A_n))$, where $\mathrm{kl}$ is the Bernoulli divergence. Write $p_n=P(A_n)$ and $q_n=Q(A_n)\le s$. Then $\mathrm{kl}(p_n,q_n)=p_n\log(1/q_n)+(1-p_n)\log\{1/(1-q_n)\}-h(p_n)\ge p_n\log(1/s)-\log2$, since the middle term is nonnegative and the binary entropy satisfies $h\le\log2$. Let $n\to\infty$: $\mathbb E_P[\tau_n]\uparrow\mathbb E_P[\tau]$ and $p_n\uparrow P(A)\ge1-\beta$. Divide, and take the infimum over $Q$.
\emph{(b)} Validity is Ville's inequality, as in Proposition~\ref{prop:valid}; the factors $1+\lambda(X_i-m_0)$ are nonnegative for every $\lambda\le1/m_0$, the boundary included. The increments $Y_i=\log(1+\lambda(X_i-m_0))$ are i.i.d.\ with $Y_i\le\log(1+\lambda(1-m_0))\le\log(1/m_0)$ and mean $g_P(\lambda)>0$; finiteness of $g_P(\lambda)$ makes them integrable even at $\lambda=1/m_0$, where they are not bounded below. Write $S_n=\sum_{i\le n}Y_i$. Wald's identity for the bounded stopping time $\tau\wedge n$ gives $g_P(\lambda)\,\mathbb E[\tau\wedge n]=\mathbb E[S_{\tau\wedge n}]$. Before $\tau$ the walk is below $\log(1/a)$, and the step that crosses adds at most $\log(1/m_0)$, so $S_{\tau\wedge n}<\log(1/a)+\log(1/m_0)$ (the overshoot bound). Hence $g_P(\lambda)\,\mathbb E[\tau\wedge n]<\log(1/a)+\log(1/m_0)$ for every $n$, and monotone convergence gives the stated bound for $\mathbb E[\tau]$. This covers the boundary case $\lambda^*=1/m_0$ of Theorem~\ref{thm:price}(b) and Eq.~\eqref{eq:sandwich}: if $\lambda^*=1/m_0$ maximizes $g_P$, then $g_P(1/m_0)=K_{\inf}$ is finite and positive, which forces $P(X=0)=0$. Power one follows from the strong law.
\emph{(c)} \emph{Weak duality.} Let $Q$ be a null law with $P\ll Q$ and let $0\le\lambda\le1/m_0$. By Jensen's inequality,
\begin{align*}
g_P(\lambda)-\mathrm{KL}(P\|Q)&=\mathbb E_P\log\Big\{(1+\lambda(X-m_0))\frac{dQ}{dP}\Big\}\le\log\mathbb E_Q[1+\lambda(X-m_0)]\\
&\le\log\{1+\lambda(\mathbb E_QX-m_0)\}\le0.
\end{align*}
\emph{Attainment.} The function $g_P$ is concave on $[0,1/m_0]$; let $\lambda^*$ maximize it. It is positive because $g_P'(0)=\mathbb E_PX-m_0>0$. Define $Q^*(x)=P(x)/(1+\lambda^*(x-m_0))$ on the support of $P$ and put the leftover mass $1-Q^*(\operatorname{supp}P)$ at $x=0$. Since $1/(1+\lambda u)=1-\lambda u/(1+\lambda u)$, the total mass of the first part is $1-\lambda^*\mathbb E_P[(X-m_0)/(1+\lambda^*(X-m_0))]=1-\lambda^*g_P'(\lambda^*)$.
\begin{itemize}
\item If $\lambda^*<1/m_0$, then $g_P'(\lambda^*)=0$. The first part has mass $1$, there is no leftover, and $\mathbb E_{Q^*}[X-m_0]=g_P'(\lambda^*)=0$.
\item If $\lambda^*=1/m_0$, then $g_P$ is finite only if $P(X=0)=0$, and $1+\lambda^*(x-m_0)=x/m_0$. Optimality at the right endpoint gives $g_P'(1/m_0)\ge0$, so the leftover mass $1-m_0\mathbb E_P[1/X]$ is nonnegative. The mean is $\mathbb E_{Q^*}[X-m_0]=(m_0-m_0^2\mathbb E_P[1/X])-m_0(1-m_0\mathbb E_P[1/X])=0$.
\end{itemize}
In both cases $Q^*$ is a null law with $\mathrm{KL}(P\|Q^*)=\mathbb E_P\log\{P(X)/Q^*(X)\}=g_P(\lambda^*)$. Combined with weak duality this gives $\max_\lambda g_P=K_{\inf}(P,m_0)$, attained at $Q^*$. Concavity makes the supremum over $[0,1/m_0)$ equal to the maximum over $[0,1/m_0]$. For laws with infinite support the same argument applies with sums replaced by integrals.

\paragraph{Theorem~\ref{thm:price}.}
\emph{(a)} Take the null law $Q^*$ of Proposition~\ref{thm:onexam}(c), for which $\mathrm{KL}(P\|Q^*)=K_{\inf}<\infty$, and let $A_k$ be the event that exam $k$ admits, so $Q^*(A_k)=\bar\alpha q_k$. Proposition~\ref{thm:lineage}(b) gives $\sum_kq_k\le1$. Follow the proof of Proposition~\ref{thm:onexam}(a) with $Q=Q^*$, but keep the binary entropy: $\mathbb E_P[\tau_{k,n}]K_{\inf}\ge\mathrm{kl}(p_n,s_n)\ge p_n\log(1/(\bar\alpha q_k))-h(p_n)$, where $p_n=P(A_k\cap\{\tau_k\le n\})$ and $s_n\le\bar\alpha q_k$. Let $n\to\infty$. Then $p_n\uparrow p=P(A_k)\ge1-\beta\ge1/2$, $h$ is continuous, and $h(p)\le h(1-\beta)=h(\beta)$ because $h$ decreases on $[1/2,1]$. So $\mathbb E_P[\tau_k]K_{\inf}\ge(1-\beta)\log\{1/(\bar\alpha q_k)\}-h(\beta)$. Exam $k$ is fixed in advance and runs on variants revealed after $\{J=k\}$ is determined, which are i.i.d.\ $P$ given that event and the earlier history, so $\mathbb E[\tau_J\mid J=k]=\mathbb E_P[\tau_k]$. Average over $k\sim\pi$, the marginal law of $J$, and write $\sum_k\pi_k\log(1/q_k)=H(\pi)+\mathrm{KL}(\pi\|q)$. Finally $\mathrm{KL}(\pi\|q)\ge-\log\sum_kq_k\ge0$ by Gibbs' inequality for a subprobability.
\emph{(b)} The gate is lineage-valid by Proposition~\ref{thm:lineage}(a), since $\sum_k\alpha q_k\le\alpha$. Proposition~\ref{thm:onexam}(b) at level $\alpha q_k$ gives $\mathbb E[\tau_k]\le\{\log(1/\alpha)+\log(1/q_k)+\log(1/m_0)\}/g_P(\lambda)$. As in part (a), $\mathbb E[\tau_J\mid J=k]=\mathbb E_P[\tau_k]$. Average over $\pi$, and use $\sum_k\pi_k\log(1/q_k)=H(\pi)+\mathrm{KL}(\pi\|q)$.

\paragraph{Eq.~\eqref{eq:sandwich}.} The right inequality is part (b) with $\lambda=\lambda^*$, since $g_P(\lambda^*)=K_{\inf}$ by Proposition~\ref{thm:onexam}(c). For the left one, the exams have power one by Proposition~\ref{thm:onexam}(b), so $p=1$ and $h(p)=0$ in the proof of part (a). Exam $k$ is valid at level $\alpha q_k$, so $Q^*(A_k)\le\alpha q_k$, and the proof of part (a) gives $\mathbb E_P[\tau_k]K_{\inf}\ge\log\{1/Q^*(A_k)\}\ge\log\{1/(\alpha q_k)\}$. Average over $\pi$.

\paragraph{Corollary~\ref{cor:universal}.}
If $\log(1/q_k)\le\log k+\log\log k$ for all large $k$, then $q_k\ge1/(k\log k)$ for all large $k$, and $\sum_kq_k=\infty$ because $\sum_k1/(k\log k)$ diverges. For the second claim, $c=\sum_k1/\{(k+1)\log^2(k+1)\}$ is finite, and $q_k=c^{-1}/\{(k+1)\log^2(k+1)\}$ gives $\log(1/q_k)=\log(k+1)+2\log\log(k+1)+\log c$.

\paragraph{Proposition~\ref{prop:minimax}.}
\emph{(a)} Let $q$ be the gate's schedule under $Q^*$, as in Theorem~\ref{thm:price}(a), with total mass $t=\sum_kq_k\le1$ and normalization $\tilde q=q/t$. Then $\mathrm{KL}(\pi\|q)=\mathrm{KL}(\pi\|\tilde q)+\log(1/t)\ge\mathrm{KL}(\pi\|\tilde q)$. Theorem~\ref{thm:price}(a) gives $K_{\inf}C(\pi)-(1-\beta)\{\log(1/\bar\alpha)+H(\pi)\}\ge(1-\beta)\mathrm{KL}(\pi\|\tilde q)-h(\beta)$, and $\sup_{\pi\in\Pi}\mathrm{KL}(\pi\|\tilde q)\ge R(\Pi)$.
\emph{(b)} Take the supremum over $\pi$ of the right inequality of Eq.~\eqref{eq:sandwich}.
\emph{(c)} For any $q$, expanding $\log(\pi_{jk}/q_k)=\log(\pi_{jk}/\bar\pi_{w,k})+\log(\bar\pi_{w,k}/q_k)$ gives
$\sup_{\pi\in\Pi}\mathrm{KL}(\pi\|q)\ge\sum_jw_j\mathrm{KL}(\pi_j\|q)=\sum_jw_j\mathrm{KL}(\pi_j\|\bar\pi_w)+\mathrm{KL}(\bar\pi_w\|q)\ge\sum_jw_j\mathrm{KL}(\pi_j\|\bar\pi_w)$.
Take the infimum over $q$.
\emph{(d)} Since $\sum_{k\le N}q_k\le1$, some $k\le N$ has $q_k\le1/N$, and the point mass at $k$ has $\mathrm{KL}=\log(1/q_k)\ge\log N$. For the uniform $u$, $\mathrm{KL}(\pi\|u)=\log N-H(\pi)\le\log N$. If $\pi$ is non-increasing, $k\pi_k\le\sum_{j\le k}\pi_j\le1$, so $\log\{\pi_k/q_k\}=\log(k\pi_kH_N)\le\log H_N$ for $q_k=1/(kH_N)$.
\emph{Certified interval in E12.} Let $U_m$ be uniform on $\{1,\dots,m\}$. Every $\pi\in\Pi_N^{\downarrow}$ is the mixture $\sum_{m\le N}m(\pi_m-\pi_{m+1})U_m$, with $\pi_{N+1}=0$; the weights are nonnegative and sum to $\sum_k\pi_k=1$. Since $\mathrm{KL}(\pi\|q)$ is convex in $\pi$, $\sup_{\pi\in\Pi_N^\downarrow}\mathrm{KL}(\pi\|q)=\max_m\mathrm{KL}(U_m\|q)$ for every $q$. Hence, for any weights $w$ and any $q$, part (c) gives $\sum_mw_m\mathrm{KL}(U_m\|\bar U_w)\le R(\Pi_N^\downarrow)\le\max_m\mathrm{KL}(U_m\|q)$. E12 evaluates both ends at the prior and output law of the Blahut-Arimoto iteration, so the interval does not depend on its convergence.

\paragraph{Theorem~\ref{thm:learned}.}
\emph{(a)} The counts, and hence $\hat q^{(i)}$, are fixed when lineage $i$ opens. They are part of the history $\mathcal H_0$ of lineage $i$, whatever produced them. Also $\sum_k\hat q^{(i)}_k=(1-\epsilon)+\epsilon=1$. The proof of Proposition~\ref{thm:lineage}(a), conditional on $\mathcal H_0$, gives $\Pr(\text{ever admitted}\mid\mathcal H_0)\le\sum_k\alpha\hat q^{(i)}_k=\alpha$.
\emph{(b)} The counted lineages are chosen without regard to their attempt indices, so the honest counts at $k$ are $b_k\sim\mathrm{Bin}(m_h,\pi_k)$. (If lineages were counted because they had closed, short lineages would be over-represented and this step would fail.) Then $c_k\ge b_k$ and $m\le m_h+r$, so $\hat q^{(i)}_k\ge(1-\epsilon)(b_k+1)/(m_h+r+N)$ on every sample. Hence
$\mathbb E\log(\pi_k/\hat q^{(i)}_k)\le\log\frac1{1-\epsilon}+\log\pi_k+\log(m_h+r+N)+\mathbb E\log\frac1{b_k+1}$.
By Jensen's inequality, $\mathbb E\log\{1/(b_k+1)\}\le\log\mathbb E\{1/(b_k+1)\}$, and for $b\sim\mathrm{Bin}(m,p)$,
$\mathbb E\{1/(b+1)\}=\frac1{(m+1)p}\sum_{c=0}^m\binom{m+1}{c+1}p^{c+1}(1-p)^{m-c}=\frac{1-(1-p)^{m+1}}{(m+1)p}\le\frac1{(m+1)p}$.
So $\mathbb E\log(\pi_k/\hat q^{(i)}_k)\le\log\frac1{1-\epsilon}+\log\frac{m_h+r+N}{m_h+1}$. Average over $k\sim\pi$. The pathwise bound follows from $\hat q^{(i)}_k\ge\epsilon/N$.
\emph{(c)} For $i>D+1$, lineage $i$ counts the $m_h=i-1-D$ lineages $1,\dots,i-1-D$. All are closed when lineage $i$ opens, by assumption, and the set is fixed by the indices alone, so their attempt indices are i.i.d.\ $\pi$ and (b) applies with $r=0$. For $i\le D+1$, bound the logarithm in (b) by $\log N$, its value at $m_h=0$. For $i>D+1$, use $\log(1+x)\le x$ and $m_h+1\ge i-D$: $\sum_{i=D+2}^n(N-1)/(i-D)\le(N-1)\sum_{j=2}^n1/j\le(N-1)\log n$. For the converse, let $J_1,\dots,J_n$ be i.i.d.\ $\pi$ and let $q^{(i)}$ depend on $J_1,\dots,J_{i-1}$; a rule with a delay is a special case. Then $M(j_1,\dots,j_n)=\prod_iq^{(i)}_{j_i}(j_1,\dots,j_{i-1})$ is a probability law on $\{1,\dots,N\}^n$, or a subprobability, which only increases the divergence. By the chain rule, $\sum_i\mathbb E\,\mathrm{KL}(\pi\|q^{(i)})=\mathrm{KL}(\pi^{\otimes n}\|M)$. Its supremum over $\pi\in\Pi_N$ is at least $R_n(N)=\inf_M\sup_\pi\mathrm{KL}(\pi^{\otimes n}\|M)$, whose expansion is that of \citet{xie1997minimax} \citep[see also][Theorem~2]{boucheron2009coding}. For the cost, apply the proof of Theorem~\ref{thm:price}(a) to lineage $i$ given $J_1,\dots,J_{i-1}$, with $\beta=0$, then sum and use the chain rule as above.

\emph{Capped exams.} Let $c_k$ count the counted lineages admitted at attempt $k$, and let $\tilde\pi$ be the law of the admitting attempt given admission. Lineages are independent and counted without regard to their attempt indices, so given their number the admitted honest counts are multinomial with probabilities $\tilde\pi$, and the proofs of (b) and (c) apply verbatim with $\tilde\pi$ for $\pi$ and $m_h$ the number of admitted honest lineages. For the total-variation bound, let $\hat J$ be the admitting attempt and $A$ the event that no useless version is admitted before attempt $J$ and the helpful exam at attempt $J\le N$ admits before the cap. Then $\Pr(A^c)\le\alpha+\beta_{\max}=\varepsilon$, by lineage control for the useless versions and the power assumption for exam $J$, and $\hat J=J$ on $A\subseteq\{\text{admitted}\}$. For any set $S$, $|\Pr(J\in S\mid A)-\Pr(J\in S)|\le\Pr(A^c)\le\varepsilon$. The law of $\hat J$ given admission mixes its law given $A$, with weight $\Pr(A)/\Pr(\text{admitted})$, and its law on the rest of the admission event, whose weight is at most $\varepsilon/(1-\varepsilon)$. Hence $\mathrm{TV}(\tilde\pi,\pi)\le\varepsilon+\varepsilon/(1-\varepsilon)$. Lineages are taken independent, as in part (b). No assumption on resubmission after a miss is needed; the cap of $N$ attempts is what makes a uniform $\beta_{\max}$ possible, since an exam truncated at $N_{\max}$ cannot admit at a level below $m_0^{N_{\max}}$.

\paragraph{Theorem~\ref{thm:adaptive}.} \emph{(a)} Since $\lambda_j\le\rho/m_0<1/m_0$, each factor $1+\lambda_j(X_i-m_0)\ge1-\rho>0$, and under Definition~\ref{def:condnull} each wealth $W_n(\lambda_j)=\prod_{i\le n}\{1+\lambda_j(X_i-m_0)\}$ is a nonnegative supermartingale with $W_0=1$. By monotone convergence so is $M_n=\sum_jw_jW_n(\lambda_j)$, with $M_0=\sum_{j\ge0}w_j=c_u\sum_{k\ge1}1/\{(k+1)\log^2(k+1)\}=1$. Ville's inequality gives $\Pr(\sup_nM_n\ge1/a)\le a$. Nothing in the exam depends on $P$. Power one follows from (b), since $\mathbb E_P[\tau]<\infty$.
\emph{(b)} The maximizer exists: $g_P$ is concave on $[0,1/m_0]$ and upper semicontinuous, since $\log\{1+\lambda(x-m_0)\}\le\log(1/m_0)$ allows the reverse Fatou lemma; its maximum is $K_{\inf}$ by Proposition~\ref{thm:onexam}(c), and $\lambda^*>0$ because $g_P'(0)=\mathbb E_PX-m_0>0$. Let $j=j^*$. Then $j\le\Lambda/\log(1/\rho)<j+1$, so $\rho^{j+1}<m_0\lambda^*\le\rho^{j}$, and $\rho\lambda^*\le\lambda_j\le\lambda^*$. Concavity with $g_P(0)=0$ gives $g_P(\lambda_j)\ge(\lambda_j/\lambda^*)g_P(\lambda^*)\ge\rho K_{\inf}$. Since $M_n\ge w_jW_n(\lambda_j)$, the exam stops no later than the single-bet exam with bet $\lambda_j$ at level $aw_j$, and Proposition~\ref{thm:onexam}(b) bounds the latter's expected stopping time by $\{\log(1/a)+\ell(j)+\log(1/m_0)\}/g_P(\lambda_j)$. Divide by $\rho K_{\inf}$, and use $1/\rho=1+1/L$. For the explicit form, $\log(1/\rho)=\log(1+1/L)\ge1/(L+1)$ gives $j^*\le(L+1)\Lambda$, so $j^*+2\le(L+1)(2+\Lambda)$. Since $(j+2)\log^2(j+2)$ increases in $j$, $\ell(j^*)=\log(1/c_u)+\log(j^*+2)+2\log\log(j^*+2)\le\log(1/c_u)+u+2\log u$. Expand $(1+1/L)L=L+1$.
\emph{(c)} Validity is Proposition~\ref{thm:lineage}(a) with (a). Apply (b) to exam $k$ with $L_k=\log(1/\alpha)+\log(1/q_k)\ge\log(1/\alpha)$, so $1+1/L_k\le1+1/\log(1/\alpha)$; this replacement multiplies a nonnegative term, since $m_0\le1$, $c_u<1$ and $u_k=\log\{(L_k+1)(2+\Lambda)\}\ge\log\{2(1+\log(1/\alpha))\}\ge1$ for $\alpha\le0.5$. Then average over $k\sim\pi$ as in Theorem~\ref{thm:price}, using $\mathbb E[\tau_J\mid J=k]=\mathbb E_P[\tau_k]$. The functions $L\mapsto\log(L+1)$ and $L\mapsto\log\log\{(L+1)(2+\Lambda)\}$ are concave and increasing (the second is the logarithm of a positive concave function), so Jensen's inequality bounds their averages by their values at $\sum_k\pi_kL_k=\bar L$. The lower bound is Theorem~\ref{thm:price}(a) with $\beta=0$.

\paragraph{Scope.} Theorem~\ref{thm:price} counts only the helpful version's exam, and assumes untruncated exams with the stated power. With $N_{\max}$ truncation and futility, $\beta>0$ enters through $(1-\beta)$, and the useless earlier versions also spend variants. The theorems treat the lineage component of the deployed level $\min(\alpha_t,\alpha\gamma^{\text{lin}}_k)$, not its interaction with e-LOND. The cost bounds use i.i.d.\ variants within an exam. The validity statements use only the conditional null.
